\section{Statement, scope, and constructive reading}

\subsection{The two main results}

The phrase ``formula by radicals'' is often stated informally.  In a
constructive setting one must say what the radical symbols mean, what data are
being supplied, and what exactly is proved.  The following definition is the
one used throughout.

\begin{definition}[Radical tower and complete radical factorization]
Let $F$ be a field of characteristic zero.  A \emph{radical tower over $F$}
is a finite chain
\[
 F=F_0\subseteq F_1\subseteq\cdots\subseteq F_r
\]
with elements $\alpha_i\in F_i$ and integers $m_i\ge 2$ such that
\[
 F_i=F_{i-1}(\alpha_i),\qquad \alpha_i^{m_i}\in F_{i-1}.
\]
A polynomial $f(T)=a_nT^n+\cdots+a_0\in F[T]$, $a_n\ne0$, has a
\emph{complete radical factorization} if some radical tower contains elements
$r_1,\ldots,r_n$ satisfying the exact identity
\[
 f(T)=a_n\prod_{j=1}^{n}(T-r_j).
\]
Multiplicity is retained.  No ordering or analytic branch convention is part
of the definition.
\end{definition}

The identity is stronger than the assertion that each displayed value happens
to be a root: it gives soundness, exhaustiveness, and multiplicity in one
finite certificate.

\begin{theorem}[Constructive formula theorem through degree four]
\label{thm:low-degree-main}
For every characteristic-zero field and every nonzero polynomial of degree at
most four, the classical linear, quadratic, Cardano, and Ferrari constructions
admit the following intuitionistically valid formulation.

Given explicit radical witnesses satisfying the equations displayed in
Sections~\ref{sec:quadratic}, \ref{sec:cubic}, and \ref{sec:quartic}, a fixed
primitive-recursive expression returns a list of the required length, and an
exact factorization identity proves that this list contains precisely all
roots with multiplicity.  Compatible radical witnesses can be organized in a
radical tower.  No proof uses excluded middle for an arbitrary proposition;
the only case distinctions are equality tests on coded field expressions or
natural numbers.
\end{theorem}

\begin{theorem}[Constructive Abel--Ruffini theorem]
\label{thm:abel-main}
There is no complete formula by radicals for the general polynomial of degree
$n$ when $n\ge5$.  More precisely, let $x_1,\ldots,x_5$ be independent
indeterminates, let
\[
 L=\Q(x_1,\ldots,x_5),
\]
let $e_i$ be the elementary symmetric functions, and let
$K=\Q(e_1,\ldots,e_5)\subseteq L$.  Then
\[
 P(T)=\prod_{i=1}^{5}(T-x_i)\in K[T]
\]
has no complete radical factorization over $K$.  Consequently a coefficient-
parametrized radical solver valid for every quintic cannot exist.  For every
$n>5$, the polynomial $T^{n-5}P(T)$ gives the same obstruction in exact degree
$n$.

After coding all finite objects by natural numbers, the proof is formalizable
in intuitionistic primitive recursive arithmetic.  In particular, ProveIt's
native concise calculus proves the translation of the arithmetized theorem by
the established bi-interpretation with iPRA.
\end{theorem}

\subsection{What Abel--Ruffini does and does not say}

The negative theorem concerns a \emph{uniform complete} radical solution for
all polynomials of a fixed degree.  It does not say that every individual
quintic is unsolvable.  For example, $T^5-1$ is already solved by roots of
unity, and a reducible quintic may reduce to equations of smaller degree.  It
also does not forbid algorithms using operations beyond radicals, nor does it
forbid numerical approximation.

The word \emph{complete} matters.  A polynomial can have one easy root and
other roots that are not obtainable by radicals.  A purported solution of the
equation must account for the whole factorization, not merely produce one
root.  The generic polynomial used here is stronger still: its Galois group is
$S_5$, and no radical tower can contain all five independent roots.

\subsection{Why radical witnesses replace branch functions}

A complex square root is not a single-valued analytic function on all of
$\C$.  A cube root has three values, and Cardano's formula fails if two cube
roots are selected independently.  Therefore the constructive statement is
relational:
\[
  s^2=D,
  \qquad u^3=R,
  \qquad v^3=S,
  \qquad uv=-p/3.
\]
The data are the elements $s,u,v$ together with proofs of these equations.
The formula then proves a polynomial identity.  This is exactly the style of
the existing ProveIt Lean development: square and cube roots are represented
by field elements plus their defining equations, and the formal theorems are
generic over characteristic-zero fields.

This witness-oriented formulation is more constructive than choosing a global
branch.  It is also stronger algebraically: it works in every field extension
where the witnesses exist, not only in $\C$ and not only for a preferred
analytic branch.

\subsection{Relationship to the current ProveIt development}

At the inspected repository revision, the following semantic Lean layers are
already present.

\begin{itemize}
\item \code{Algebra/PolynomialFormulas/Lean/PolynomialFormulas/Basic.lean}
  proves the linear and quadratic formulae, exact quadratic factorization, and
  quadratic exhaustiveness.
\item \code{Cubic.lean} proves the depressed-cubic transformation, compatible
  Cardano radicals, the three-factor identity, and
  \code{cubic_eq_zero_iff}.
\item \code{Quartic.lean} proves Ferrari's resolvent equations, the two-
  quadratic factorization, and \code{quartic_eq_zero_iff}.
\item \code{GenericAbelRuffini.lean} constructs a generic polynomial whose
  Galois group is a symmetric group and derives a rootwise radical
  obstruction through mathlib's field-theoretic infrastructure.
\item \code{Logic/PrimitiveRecursiveArithmetic/ConciseIPRA} supplies the
  formal iPRA presentation and the theorem \code{native_iff_pra} transferring
  derivability between the native calculus and ordinary iPRA.
\end{itemize}

The present report does not replace those files.  It reorganizes their
mathematics into a deliberately finite, certificate-oriented proof whose
arithmetization uses only primitive recursion and quantifier-free induction.
Section~\ref{sec:lean-plan} identifies the exact additional bridge modules
needed for a fully kernel-checked iPRA result.

\subsection{Constructive discipline used below}

The proof observes four rules.

\begin{enumerate}
\item A radical is always accompanied by its power equation.
\item Division is performed only with an explicit nonzero hypothesis.
\item A disjunction is introduced only after a decidable finite equality test,
  or it is avoided by retaining a factorization identity.
\item Every finite construction is given with a primitive-recursive bound;
  no appeal is made to unbounded minimization, a global choice function, or a
  completed infinite set.
\end{enumerate}

These rules will be repeated at the few places where a classical textbook
proof usually suppresses a branch or an existence argument.
